\begin{lemma}[The resolved leg is embedded]\label{lem:embedded-leg}
The second projection
\begin{equation}\label{eq:e-embedding}
  e\colon P_K\hookrightarrow\widehat X_t
\end{equation}
is a topological embedding.
\end{lemma}

\begin{proof}
First consider the second projection on the entire pullback \( P \).  Its image is \( p_t^{-1}(a(X_s)) \).  Because \( a \) is an embedding, the formula
\begin{equation}\label{eq:e-inverse}
  \widehat y\longmapsto\bigl(a^{-1}(p_t(\widehat y)),\,\widehat y\bigr)
\end{equation}
is a continuous inverse on this image.  Thus the second projection identifies \( P \) homeomorphically with \( p_t^{-1}(a(X_s)) \).  Restricting this homeomorphism to \( P_K \) proves \eqref{eq:e-embedding}.
\end{proof}

\begin{remark}
The embedding hypothesis on \( a \) is essential: if \( a \) has a
base-space self-intersection, equality of second coordinates in
\eqref{eq:pullback} need not determine the first coordinate, and
Lemma~\ref{lem:embedded-leg} can fail before any subgroup collision is
considered; compare Section~\ref{sec:seedvsbold}.
\end{remark}